\documentclass{modern-c4}

\title{Python OOP Guide for Data Scientists: Building Production-Ready ML Systems}
\author{Engineering Team}
\date{\today}

\begin{document}

\maketitle

\begin{abstract}
This guide bridges the gap between experimental Jupyter Notebooks and production-grade Object-Oriented Python. It outlines a structured, orthogonal architecture that isolates state, embraces SOLID principles, and provides a scalable foundation for Machine Learning projects.
\end{abstract}

\section{1. The Motivation: From Notebooks to Engineering}
Jupyter Notebooks are unparalleled for Exploratory Data Analysis (EDA) and rapid prototyping. However, migrating notebook scripts directly to production often leads to:
\begin{itemize}
    \item \textbf{Hidden State Dependencies:} Out-of-order cell execution causing unreproducible results or global variable contamination.
    \item \textbf{Spaghetti Code:} Data loading, feature engineering, and training logic tightly coupled in long linear scripts.
    \item \textbf{Data Leakage:} Transforming testing data using testing statistics rather than isolating training statistics cleanly.
\end{itemize}
Our standardized Python OOP template solves these issues by establishing a clear "Left-Center-Right" modular architecture.

\section{2. High-Level Conceptual Architecture}
The system architecture visually and logically separates concerns. We enforce a strict boundary between user interaction, the ML core logic, and external data sources.

\begin{center}
\begin{tikzpicture}[node distance=2.5cm, auto]
    % --- LEFT ZONE: User / Context ---
    \node[c4_person] (datascientist) {
        Data Scientist \\
        \normalfont\scriptsize Executes experiments \& pipelines.
    };

    \node[c4_external, below=1cm of datascientist] (notebooks) {
        Jupyter Notebooks \\
        \normalfont\scriptsize EDA \& Prototyping
    };

    % --- CENTER ZONE: System Core ---
    \node[c4_container, right=3.5cm of datascientist, yshift=2.5cm] (pipeline) {
        Training Pipeline \\
        \normalfont\scriptsize Orchestrates the ML Workflow
    };

    \node[c4_container, below=0.6cm of pipeline] (preprocessor) {
        Preprocessor \\
        \normalfont\scriptsize Stateful Feature Engineering
    };

    \node[c4_container, below=0.6cm of preprocessor] (trainer) {
        Model Trainer \\
        \normalfont\scriptsize Algorithm Encapsulation
    };

    % Draw a bounding box for the system core
    \begin{scope}[on background layer]
        \node[c4_system, fit=(pipeline) (preprocessor) (trainer), inner sep=15pt, label={[font=\sffamily\bfseries\color{SystemBlue}]above:System Boundary (ML Template)}] (system_boundary) {};
    \end{scope}

    % --- RIGHT ZONE: Adapters & Resources ---
    \node[c4_system, fill=DbGreen!10, draw=DbGreen, right=1.5cm of pipeline] (adapters) {
        Adapters (loader.py) \\
        \normalfont\scriptsize External I/O Interfaces
    };

    \node[c4_database, below=1.5cm of adapters] (db) {
        Data Lake / DB \\
        \normalfont\scriptsize Raw \& Processed Data
    };

    % --- Connections ---
    \draw[c4_rel] (datascientist) -- node[above, sloped, pos=0.55] {Triggers CLI} (system_boundary.west |- pipeline);
    \draw[c4_rel, dashed] (datascientist) -- node[left] {Explores in} (notebooks);

    \draw[c4_rel] (pipeline) -- node[right] {Uses} (preprocessor);
    \draw[c4_rel] (pipeline) to[bend right=45] node[left] {Uses} (trainer);

    \draw[c4_rel] (pipeline) -- node[above, pos=0.65] {Requests Data via} (adapters);
    \draw[c4_rel] (adapters) -- node[right] {Queries} (db);

\end{tikzpicture}
\end{center}

\section{3. Architectural Principles}
\subsection{Orthogonal Dimensions}
To prevent state contamination, the architecture is strictly divided into three distinct dimensions:
\begin{enumerate}
    \item \textbf{Control Dimension (\texttt{main.py}, \texttt{pipeline.py}):} Orchestrates the workflow. It defines the sequence of execution without hardcoding any mathematical operations or rules.
    \item \textbf{Domain Dimension (\texttt{preprocessor.py}, \texttt{trainer.py}):} Encapsulates the business logic and ML state (e.g., standard deviation for scaling, tree weights, one-hot encoders).
    \item \textbf{Cross-Cutting Dimension (\texttt{config.py}, \texttt{utils/}):} Infrastructure that spans the entire system, such as pure utility functions and the Pydantic-based configuration (12-Factor App approach).
\end{enumerate}

\subsection{Dependency Inversion}
We adhere to SOLID principles, primarily Dependency Inversion (DI). The \texttt{TrainingPipeline} does not instantiate a concrete \texttt{DataLoader} directly. Instead, concrete implementations are injected via constructors in the composition root (\texttt{main.py}). This allows components to be independently testable using mocks without triggering real network requests.

\section{4. Core Component Breakdown}
\begin{itemize}
    \item \textbf{\texttt{adapters/}:} Houses logic to communicate with external systems. It abstracts away SQL queries or CSV reading into generic protocols. This protects the ML Core from database connection failures.
    \item \textbf{\texttt{ml\_core/}:} The heart of the system. The \texttt{Preprocessor} utilizes \texttt{fit\_transform()} to learn state and \texttt{transform()} for inference. This ensures that validation data is processed solely using training statistics, strictly preventing data leakage.
    \item \textbf{\texttt{utils/}:} Pure, stateless mathematical or statistical functions that can be safely imported by both notebooks and production pipelines to guarantee identical computational behavior across environments.
\end{itemize}

\end{document}
